\begin{tikzpicture}[
  font=\scriptsize,
  sel/.style={draw, rounded corners, minimum height=1.05cm, minimum width=2.7cm,
              text width=2.5cm, align=center, font=\scriptsize\bfseries},
  label/.style={draw, rounded corners, fill=praditagreen!15, align=center,
                font=\scriptsize\bfseries},
  panah/.style={-{Stealth[length=2mm]}, thick}
]
  % Judul kolom: tiga aspek inti, ditambah aspek motivasi di kanan.
  \node[label, minimum width=2.7cm, minimum height=0.6cm] at (2.95,3.55) {Struktur pasif};
  \node[label, minimum width=2.7cm, minimum height=0.6cm] at (5.75,3.55) {Perilaku};
  \node[label, minimum width=2.7cm, minimum height=0.6cm] at (8.55,3.55) {Struktur aktif};

  % Lapisan strategi di atas inti.
  \node[sel, fill=blue!8, minimum width=8.3cm, text width=8.1cm, minimum height=0.75cm]
    at (5.75,2.6)
    {Strategi: kapabilitas, \textit{value stream}, sumber daya, \textit{course of action}};

  % Judul baris: tiga lapisan inti.
  \node[label, minimum width=1.5cm, minimum height=1.05cm] at (0.75,1.5) {Bisnis};
  \node[label, minimum width=1.5cm, minimum height=1.05cm] at (0.75,0.3) {Aplikasi};
  \node[label, minimum width=1.5cm, minimum height=1.05cm] at (0.75,-0.9) {Teknologi};

  % Isi inti 3 x 3, dengan contoh dari layanan janji temu Medikara.
  \node[sel, fill=orange!12] at (2.95,1.5)
    {Objek bisnis\\{\tiny\mdseries Janji Temu}};
  \node[sel, fill=orange!12] at (5.75,1.5)
    {Proses, fungsi, layanan bisnis\\{\tiny\mdseries Membuat Janji Temu}};
  \node[sel, fill=orange!12] at (8.55,1.5)
    {Aktor, peran\\{\tiny\mdseries Pasien, Petugas Pendaftaran}};

  \node[sel, fill=orange!12] at (2.95,0.3)
    {Objek data\\{\tiny\mdseries Data Jadwal Dokter}};
  \node[sel, fill=orange!12] at (5.75,0.3)
    {Fungsi, layanan aplikasi\\{\tiny\mdseries Layanan Pemesanan Jadwal}};
  \node[sel, fill=orange!12] at (8.55,0.3)
    {Komponen, antarmuka\\{\tiny\mdseries Medikara+, SIMRS}};

  \node[sel, fill=orange!12] at (2.95,-0.9)
    {Artefak\\{\tiny\mdseries berkas basis data}};
  \node[sel, fill=orange!12] at (5.75,-0.9)
    {Layanan teknologi\\{\tiny\mdseries Layanan Basis Data}};
  \node[sel, fill=orange!12] at (8.55,-0.9)
    {\textit{Node}, perangkat\\{\tiny\mdseries Server Basis Data SIMRS}};

  % Aspek motivasi di sisi kanan, merentang seluruh lapisan.
  \node[sel, fill=blue!8, minimum width=1.8cm, text width=1.6cm, minimum height=4.35cm]
    at (11.0,0.3)
    {Motivasi\\[3pt]{\tiny\mdseries \textit{stakeholder}, \textit{driver}, \textit{goal},
     \textit{outcome}, prinsip, \textit{requirement}}};

  % Lapisan implementasi dan migrasi di bawah inti.
  \node[sel, fill=blue!8, minimum width=8.3cm, text width=8.1cm, minimum height=0.75cm]
    at (5.75,-2.0)
    {Implementasi dan migrasi: \textit{work package}, \textit{deliverable}, \textit{plateau}, \textit{gap}};

  \node[font=\tiny, gray, anchor=north] at (5.75,-2.45)
    {kotak oranye: kerangka inti tiga lapisan kali tiga aspek};
\end{tikzpicture}
